\chapter{Kidney Pain}

\section*{Definition}
Pain originating from the kidneys, typically felt in the flank region (costovertebral angle) between the ribs and hip, often radiating to the lower abdomen or groin.

\section*{Pathophysiology}
Kidney pain arises from distention of the renal capsule (stretching from inflammation, obstruction, or mass), ureteral colic (smooth muscle spasm from stone passage), or inflammation of renal parenchyma (pyelonephritis). Unlike musculoskeletal back pain, renal pain is typically constant, not positional, and associated with urinary symptoms.

\section*{Etiology}
\begin{itemize}
  \item Nephrolithiasis (kidney stones)
  \item Pyelonephritis (kidney infection)
  \item Ureteral obstruction
  \item Renal infarction or thrombosis
  \item Polycystic kidney disease
  \item Renal cell carcinoma
  \item Glomerulonephritis
  \item Hydronephrosis
  \item Papillary necrosis
  \item Perinephric abscess
\end{itemize}

\section*{Indications for Evaluation}
Seek care if:
\begin{itemize}
  \item Severe or worsening symptoms
  \item Severe flank pain with fever and chills, inability to urinate, blood in urine, or severe unrelenting pain
  \item Accompanied by other concerning signs
\end{itemize}

\section*{Related Symptoms}
\begin{itemize}
  \item Flank pain
  \item Blood in urine (hematuria)
  \item Fever
  \item Nausea and vomiting
  \item Dysuria
\end{itemize}
\textbf{Note:} Always consider serious underlying conditions when evaluating persistent kidney pain.

\section*{Self-Care / Home Management}
\begin{itemize}
  \item Rest and avoid activities that may aggravate the condition.
  \item Maintain adequate hydration and a balanced diet.
  \item Over-the-counter remedies may provide symptomatic relief (consult a pharmacist or doctor).
  \item Monitor symptoms and seek professional advice if they persist or worsen.
  \item Avoid self-medicating with prescription medications without medical guidance.
\end{itemize}

\section*{Diagnostic Approach}
\begin{itemize}
  \item A thorough history and physical examination are the first steps in evaluation.
  \item Laboratory tests (blood work, urinalysis) may be ordered based on clinical suspicion.
  \item Imaging studies (X-ray, ultrasound, CT, MRI) may be indicated when structural pathology is suspected.
  \item Specialist referral is arranged when the diagnosis remains uncertain or requires specific expertise.
\end{itemize}

\section*{Treatment / Management Overview}
\begin{itemize}
  \item Treatment targets the underlying cause identified through diagnostic evaluation.
  \item Supportive care and symptom management are provided while awaiting definitive therapy.
  \item Medications, lifestyle modifications, and physical therapies may be prescribed as appropriate.
  \item Follow-up ensures treatment effectiveness and allows timely adjustments to the care plan.
\end{itemize}

\clearpage